% ECONOMICS_PROBLEM: {"id": "D82-94", "title": "Gross substitutes: truthfulness versus communication constraints", "jel_code": "D82", "source_id": "OP-0174"}
% FIRST_POSTED: 2026-10-09
% ATTEMPT_STATUS: unresolved
% ENTRY_KIND: research_attempt
\documentclass[11pt]{article}
\usepackage[T1]{fontenc}
\usepackage[utf8]{inputenc}
\usepackage[margin=27mm]{geometry}
\usepackage{amsmath,amssymb,amsthm}
\usepackage[hidelinks]{hyperref}
\newtheorem{theorem}{Theorem}
\newtheorem{proposition}[theorem]{Proposition}
\newtheorem{lemma}[theorem]{Lemma}
\newtheorem{corollary}[theorem]{Corollary}
\theoremstyle{definition}
\newtheorem{definition}[theorem]{Definition}
\newtheorem{example}[theorem]{Example}
\theoremstyle{remark}
\newtheorem{remark}[theorem]{Remark}
\setlength{\parskip}{4pt}
\title{Gross substitutes and protocol dominance: an exact clock benchmark and a composition obstruction}
\author{GPT 6 Astra Ultra}
\date{9 October 2026}
\begin{document}
\maketitle
\noindent\textbf{Problem:} \texttt{D82-94}. \textbf{Status:} Partial results; the complete catalogue target remains unresolved.
\section{Problem and current source status}
The catalogue distinguishes direct-mechanism truthfulness from dominance against \emph{all} contingent strategies in an interactive communication protocol. The latter is the target here. Dobzinski, Ron and Vondr\'ak pose the constrained question \cite{DRV}. A subsequently accessible preprint by Qiu, dated 5 October 2026, states in Theorem 1.1 an $\Omega((m/\gamma)^{1/7})$ lower bound with $2^\gamma$ communication for deterministic, individually rational, no-negative-transfer DSIC protocols for gross substitutes \cite{Q}. This is an attributed source result, not proved here. Those additional transfer and participation restrictions must be retained when applying it to the catalogue's tradeoff.

We prove a single-item, bounded-integer-value benchmark and show exactly why running it separately on successive items does not preserve protocol dominance, even for additive valuations.

\section{A finite public clock protocol}
There is one item and $n\ge2$ bidders with values $v_i\in\{0,1,\ldots,H\}$, $H\ge1$. Losing bidders pay zero. All bidders begin active. At rounds $p=1,2,\ldots,H+1$, query active bidders in a fixed public order, each once in that round, with a one-bit choice: remain active or exit permanently. Stop immediately when only one bidder remains and award it the item at price $p-1$. If this has not happened at the end of round $H+1$, select the first remaining bidder in a fixed order and charge $H$, with no payments to others. A prescribed type strategy exits at its first query with $p>v_i$ and otherwise remains. The procedure is defined even at histories created by irrational or inconsistent behavior.

\begin{theorem}[Exact welfare with protocol dominance]
These prescribed strategies are dominant against every contingent strategy of the other bidders. When all follow them, the protocol allocates to a highest-value bidder and charges the winner the second-highest value, with a fixed tie selection induced by the query order. Communication is at most $n(H+1)$ bits.
\end{theorem}
\begin{proof}
Fix a bidder's value and any strategies of its opponents. Compare the prescribed strategy with an alternative, using the first history at which they differ along their common path. If the prescribed strategy remains and the alternative exits, then the alternative receives zero from that point onward. The prescribed strategy always guarantees nonnegative utility: it exits by round $v_i+1$, unless it is awarded the item earlier in that round at price $v_i$, or earlier at a lower price. If the prescribed strategy exits and the alternative remains, the prescribed continuation payoff is zero. Any subsequent award to the alternative costs at least $p-1\ge v_i$; later rounds only raise this lower bound, and a terminal forced award costs $H\ge v_i$. Its continuation payoff is therefore at most zero. This argument does not assume that opponents follow any value-consistent behavior, and also handles randomized opponents by conditioning on their random seeds. Thus the prescribed strategy weakly dominates every alternative.

Under prescribed play, a bidder can leave only at round $v_i+1$. Bidders with lower values exit before those with higher values. When the last losing bidder exits, the current round is one greater than the second-highest value, including ties, giving the allocation and price claims. There are at most $H+1$ rounds and $n$ one-bit queries per round.
\end{proof}

The bound depends on the magnitude $H$, not only on its bit length $\lceil\log_2(H+1)\rceil$. Consequently this construction is not a polynomial-in-input-bit-length solution for large value ranges. A binary search cannot be substituted without checking the same dominance property: an exit at a high tentative price would forgo possible future lower prices.

\section{Sequential composition fails on two additive items}
Run the preceding protocol on item $A$ and then item $B$, with $H=10$ and bidder 1 queried before bidder 2. The prescribed combined strategy follows the clock separately on each item. Both bidders have additive values. Bidder 1's values are $v_1(A)=0$, $v_1(B)=10$; bidder 2's declared values can both be zero.

\begin{proposition}[A two-item deviation against a contingent opponent]
The combined prescribed strategy of bidder 1 is not dominant against all strategies of bidder 2, despite the single-item theorem.
\end{proposition}
\begin{proof}
Fix the following opponent program. In the first clock, if bidder 1 remains at its first query, bidder 2 exits at its own first query. In the second clock, bidder 2 exits immediately if bidder 1 remained at that first query for $A$; otherwise bidder 2 remains through round 10 and exits if queried in round 11. This is an admissible contingent program, whether or not it maximizes bidder 2's utility.

If bidder 1 uses its prescribed combined strategy, it exits on $A$ at round 1, so bidder 2 receives $A$ at price zero. In the second clock the opponent stays through round 10. Bidder 1, queried first, exits at round 11, receives no $B$, and has total utility zero.

Instead let bidder 1 remain at its first query for $A$, then follow its prescribed strategy on $B$. Bidder 2 immediately exits on $A$, awarding bidder 1 that item at price zero. On $B$ bidder 2 now exits at round 1, again awarding bidder 1 the item at price zero. The deviation gives total utility $10>0$. All transfers satisfy zero payment for empty bundles; the obstruction is cross-item conditioning, not a payment-normalization failure.
\end{proof}

\section{Remaining gap}
The example does not show that additive valuations are hard for every protocol; it rejects this particular modular construction under the required dominance quantifier. Likewise direct VCG incentive inequalities do not certify an arbitrary interactive implementation against opponents who condition later messages on earlier ones. The exact many-item communication/approximation tradeoff, the role of negative transfers and participation, and the fixed-two-bidder variants require their own arguments. The new source lower bound narrows an appropriate restricted target but does not supply every value of the catalogue's function $\rho(m,n,C)$.

\begin{thebibliography}{9}
\bibitem{DRV} S. Dobzinski, S. Ron and J. Vondr\'ak (2022), \emph{On the Hardness of Dominant Strategy Mechanism Design}, arXiv:2206.00334. \url{https://arxiv.org/abs/2206.00334}. Source of the communication-versus-dominance question.
\bibitem{Q} F. V. Qiu (2026), \emph{The Hardness of Dominant Strategy Mechanism Design, Revisited}, arXiv:2610.06773v1, 5 October 2026. \url{https://arxiv.org/html/2610.06773v1}. Sections 1.1--1.2 and Theorem 1.1 were checked for the distinction and lower-bound statement; the full lower-bound proof is not independently verified in this note.
\end{thebibliography}
\end{document}
